%	현대물리실험 보고서
%	실험3
%	202100973 이승엽

%----------------------------------------------------------------------------------------
%	PACKAGES AND DOCUMENT CONFIGURATIONS
%----------------------------------------------------------------------------------------

\documentclass[a4paper, 10pt, nanum]{CSUniSchoolLabReport}
% use UTF8 encoding
\usepackage[utf8]{inputenc}
% use KoTeX package for Korean
\usepackage{kotex}
\usepackage{amsmath}
\usepackage{hyperref}
\usepackage{graphicx}
\usepackage{indentfirst}
\usepackage{setspace}
\usepackage{multirow}
\usepackage{enumitem}
\usepackage{graphicx}
\usepackage{wrapfig}
\usepackage{epstopdf}
\usepackage{caption}
\usepackage{subcaption}

\setlength{\parindent}{0.2in} % 들여쓰기 길이 설정
\setlength{\parskip}{3mm} % 문단 간의 간격 조절
\setstretch{1.5} % 줄간격

\addbibresource{reference.bib} % Bibliography file (located in the same folder as the template)

%----------------------------------------------------------------------------------------
%	REPORT INFORMATION
%----------------------------------------------------------------------------------------

\title{현대물리실험 실험3 보고서 \\ Franck-Hertz experiment with neon} % Report title

\author{\textsc{Department of Physics} 202100973 이승엽}

\date{\today}

%----------------------------------------------------------------------------------------

\begin{document}

\maketitle % Insert the title, author and date using the information specified above

\begin{center}
	\begin{tabular}{l r}
		Date Performed: & April 11, 2023 \\ % Date the experiment was performed
		& April 18, 2023 \\
		Partners: & 202100969 이규리 \\ % Partner names
		& 202100989 한누리 \\
		Instructor: & Professor 이기주 \\ % Instructor/supervisor
		Typesetting: & LaTeX \\
		Datafitting: & Python \\
	\end{tabular}
\end{center}

%----------------------------------------------------------------------------------------
%	ABSTRACT
%----------------------------------------------------------------------------------------

\maketitle
% \begin{abstract}
% 	This report ...
% \end{abstract}

%----------------------------------------------------------------------------------------
%	INTRODUCTION
%----------------------------------------------------------------------------------------

\section{서론}

  To record a Franck-Hertz curve for neon. To measure the discontinuous energy emission of free electrons for inelastic collision. To interpret the measurement results as representing discrete energy absorption by neon atoms. To observe the Ne-spectral lines resulting from the electron-collision excitation of neon atoms. To identify the luminance phenomenon as layers with a high probability of excitation. 

%----------------------------------------------------------------------------------------
%	THEORY
%----------------------------------------------------------------------------------------

\section{이론}

  J. Frank와 G. Hertz가 1914년 이후 원자의 공명
퍼텐셜을 구하기 위해 실시한 실험으로서, 수은이 네
온으로 대체된 것만 빼면 같은 실험으로 당시 보어의
수소 원자 모형이 가진 전자의 Life time 문제를 설
명하는 에너지 준위의 양자화를 직접적으로 보여준
실험이다.

%----------------------------------------------------------------------------------------
%	EXPERIMENTAL METHOD
%----------------------------------------------------------------------------------------

\section{실험 방법}

	ㅇ

%----------------------------------------------------------------------------------------
%	RESULT AND DISCUSSION
%----------------------------------------------------------------------------------------

\section{실험 결과}

	ㅇ


%----------------------------------------------------------------------------------------
%	CONCLUSIONS
%----------------------------------------------------------------------------------------

\section{결론}

	ㅇ

%----------------------------------------------------------------------------------------
%	BIBLIOGRAPHY
%----------------------------------------------------------------------------------------

% \printbibliography % Output the bibliography

%----------------------------------------------------------------------------------------

\end{document}